\documentclass[oneside,a4paper,14pt]{extarticle}
\usepackage[T1,T2A]{fontenc}
\usepackage[a4paper,top=20mm,bottom=20mm,left=30mm,right=10mm]{geometry}
\usepackage{amsmath, amssymb}
\usepackage[utf8]{inputenc}
\usepackage[russian]{babel}
\usepackage{textcomp}
\usepackage{indentfirst}
\usepackage{graphicx}
\usepackage{float}
\usepackage{mwe}
\usepackage{wrapfig}
\usepackage{caption}
\usepackage{tocloft}
\renewcommand{\cftsecleader}{\cftdotfill{\cftdotsep}}
\usepackage{longtable}
\usepackage{amsfonts}
\usepackage{amsthm}
\usepackage[all]{xy}
\usepackage[breaklinks]{hyperref}
\usepackage{titlesec}
\usepackage{fancyvrb}
\usepackage{paracol}
\usepackage{listings}
\usepackage{xcolor}
\begin{document}

\newpage\thispagestyle{empty}
\begin{center}
    МИНИСТЕРСТВО НАУКИ И ВЫСШЕГО ОБРАЗОВАНИЯ\\
    РОССИЙСКОЙ ФЕДЕРАЦИИ\\
    ФЕДЕРАЛЬНОЕ ГОСУДАРСТВЕННОЕ БЮДЖЕТНОЕ\\
    ОБРАЗОВАТЕЛЬНОЕ УЧРЕЖДЕНИЕ ВЫСШЕГО ОБРАЗОВАНИЯ\\
    «ВЯТСКИЙ ГОСУДАРСТВЕННЫЙ УНИВЕРСИТЕТ»\\
    Институт математики и информационных систем\\
    Факультет автоматики и вычислительной техники\\
    Кафедра электронных вычислительных машин
\end{center}

\vspace{20mm}
\begin{center}

    Отчёт по лабораторной работе №2\\
    по дисциплине\\
    «Вычислительная математика и методы оптимизации»\\
\end{center}
\vfill

Выполнил студент гр. ИВТб-2301-05-00 \hfill
\makebox[8cm]{\hrulefill /Ковальский И.И./}

Проверил преподаватель\hfill
\makebox[8cm]{\hrulefill /Коржавина А.С./}

\vfill
\begin{center}
    Киров\\
    2026
\end{center}

\newpage
\setcounter{page}{1}
\tableofcontents
\newpage

\section{Задание}
\begin{enumerate}
    \item Изучить и реализовать заданный метод. Язык программирования выбрать самостоятельно.
    \item Продемонстрировать работу приложения на примерах. При демонстрации работы приложения выводить каждую итерацию метода (информацию о решаемых подзадачах, ветвлениях и отсечениях) отдельно, также выводить решение (ответ).
\end{enumerate}
\textbf{Заданный метод} --- метод ветвей и границ.

\section{Листинг программы}





\section{Результат работы}
\subsection{Пример 1: Простая задача}
\textbf{Задача:} $F = 3x_1 + 5x_2 \to \max$, при ограничениях $x_1 \le 4$, $x_2 \le 6$, $3x_1 + 2x_2 \le 18$.
\begin{verbatim}
--- УЗЕЛ 0 ---
LP-решение: F = 36.0000; x1=2.0000, x2=6.0000
НОВОЕ ЛУЧШЕЕ РЕШЕНИЕ! F = 36.0000

РЕЗУЛЬТАТ:
ОПТИМАЛЬНОЕ ЦЕЛОЧИСЛЕННОЕ РЕШЕНИЕ:
   x1 = 2
   x2 = 6
   Максимум F = 36.0000
\end{verbatim}

\textbf{Дерево решений:}
\begin{itemize}
    \item Узел 0 [ЦЕЛОЕ, F=36.00] (ОПТИМУМ)
\end{itemize}

\subsection{Пример 2: Задача, требующая ветвления}
\textbf{Задача:} $F = 5x_1 + 8x_2 \to \max$, при ограничениях $x_1 + x_2 \le 6$, $5x_1 + 9x_2 \le 45$.
\begin{verbatim}
--- УЗЕЛ 0 ---
LP-решение: F = 41.25, x1=2.25, x2=3.75
ВЕТВЛЕНИЕ по x2
... (промежуточные узлы) ...
--- УЗЕЛ 6 ---
LP-решение: F = 40.00, x1=0.00, x2=5.00
НОВОЕ ЛУЧШЕЕ РЕШЕНИЕ! F = 40.00
... (отсечения остальных ветвей) ...

РЕЗУЛЬТАТ:
ОПТИМАЛЬНОЕ ЦЕЛОЧИСЛЕННОЕ РЕШЕНИЕ:
   x1 = 0
   x2 = 5
   Максимум F = 40.0000
\end{verbatim}

\textbf{Дерево решений:}
\begin{itemize}
    \item Узел 0 (LP=41.25)
    \begin{itemize}
        \item Узел 1: $x_2 \le 3$ (ОТСЕЧЕН)
        \item Узел 2: $x_2 \ge 4$ (LP=41.00)
        \begin{itemize}
            \item Узел 3: $x_1 \le 1$ (LP=40.56)
            \begin{itemize}
                \item Узел 5: $x_2 \le 4$ (ОТСЕЧЕН)
                \item Узел 6: $x_2 \ge 5$ \textbf{(ЦЕЛОЕ, F=40.00 - ОПТИМУМ)}
            \end{itemize}
            \item Узел 4: $x_1 \ge 2$ (НЕТ РЕШЕНИЙ)
        \end{itemize}
    \end{itemize}
\end{itemize}

\section{Вывод}
В рамках данной лабораторной работы был реализован метод ветвей и границ для решения задач целочисленного линейного программирования. Алгоритм использует двойственный симплекс-метод, разработанный в предыдущей работе, для решения релаксированных ЛП-задач в узлах дерева поиска.

Программа успешно продемонстрировала ключевые этапы метода:
\begin{itemize}
    \item \textbf{Ветвление:} Создание подзадач путем добавления новых ограничений на дробную переменную.
    \item \textbf{Оценка:} Решение ЛП-релаксации для получения верхней оценки целевой функции.
    \item \textbf{Отсечение:} Исключение из рассмотрения ветвей, которые заведомо не могут содержать оптимального решения.
\end{itemize}
Реализация была протестирована на различных типах задач, показав свою работоспособность. Задача лабораторной работы выполнена.

\end{document}